\SourceLeaf{186}\clearpage
\begin{center}\textsc{Gauss} an \textsc{Gerling}. Göttingen, 25. August 1818.\end{center}

.\ .\ .\ .\ .\ .\ Von Ihren Fragen wegen des \textsc{Lorenz}ischen Lehrbuchs kann ich
heute nur auf einige antworten. .\ .\ .\ .\ .\ .

.\ .\ .\ .\ .\ .\ Ich freue mich, dass Sie den Muth haben sich so auszudrücken, als
wenn Sie die Möglichkeit, dass unsere Parallelentheorie, mithin unsere ganze
Geometrie, falsch wäre, anerkennten. Aber die Wespen, deren Nest Sie
aufstören, werden Ihnen um den Kopf fliegen. .\ .\ .\ .\ .\ .

\begin{center}\rule{3cm}{0.4pt}\end{center}

\begin{center}[\textsc{Gerling} \emph{an} \textsc{Gauss}. \emph{Marburg, 25. Januar 1819.}]\end{center}

\{.\ .\ .\ .\ .\ .\ .\ .\ Die Stelle über die Parallelentheorie habe ich nun so gefasst:
»Der Satz \S\ 72 ist in Euklids Elementen (1. Buch, 11. Grundsatz) als
Grundsatz aufgestellt. Dass er aber kein Grundsatz sei, sondern eines
Beweises bedürfe, lehrt die Betrachtung: dass zwei wesentlich verschiedene
Anschauungen (Winkel zweier Linien mit einer dritten und Zusammentreffen
derselben unter sich) darin vorkommen, deren \textls{nothwendiger} Zusammenhang durch
einen Beweis nachgewiesen werden muss. --- Dieser Beweis (die
Parallelentheorie) ist auf mannichfaltige Weise von scharfsinnigen
Mathematikern versucht, bis jetzt aber noch nicht vollkommen genügend
aufgefunden worden. Solange er fehlt, bleibt der Satz, so wie alles was sich
auf ihn stützt, eine Hypothese, deren Gültigkeit für unser Leben freilich
durch die Erfahrung dargethan wird, deren \textls{allgemeine, nothwendige}
Richtigkeit aber ohne Absurdität bezweifelt werden könnte«.

Ad vocem Parallelentheorie muss ich Ihnen noch etwas erzählen, und eines
Auftrags mich entledigen. Ich erfuhr im vorigen Jahr, dass mein College
\textsc{Schweikart} (Prof. juris, jetzt Prorector) sich ehedem mit
Mathematik viel beschäftigt und namentlich auch über Parallelen geschrieben
habe. Ich bat ihn also mir sein Buch zu leihen. Indem er mir diess
versprach, sagte er mir, dass er jetzt wohl einsehe, wie in seinem Buche
(1808) Fehler vorgekommen (er hatte z.\,B. Vierecke mit gleichen Winkeln als
einen ursprünglichen Begriff gebraucht), dass er aber nicht abgelassen habe,
sich mit dem Gegenstande zu beschäftigen,
